\section{DRP compute}\label{sec:upcompute}

This replaces the DRP rows of \tabref{tab:drpAndAlertSizing} and
\tabref{tab:computeSizingOps}.

\subsection{Basis}

Compute is anchored on the DR1 estimate in \emph{DRP Resource Usage
Estimates}, \S1.2.4 (\secref{sec:upmodel}):
\begin{itemize}
\item runtimes from the \vthirty{} processing (DM-53697), on Milano
      nodes;
\item input: the first year of the \texttt{baseline\_v5.0.0} simulated
      survey, \smDRoneVisits{} visits;
\item result: \smDRoneNodeDays{} node-days at \smCoresNode{} cores per node,
      or \smDRoneCoreHoursM{} million core-hours.
\end{itemize}
This supersedes the \wfortyone{} estimate of about 5{,}700 node-days (\S1.2.3).
The page attributes the increase to four or five tasks, two of which were
absent from \wfortyone{}. The figure is task-level CPU time. It excludes idle
time between tasks, gaps between pipeline stages, and calendar stretch.

\subsection{From the estimate to each operations year}

The estimate is turned into a rate per terabyte of raw input:
\[
  \text{core-hours per input TB} =
  \frac{\text{estimate core-hours}}
       {\text{estimate visits} \times \text{image size} \times \text{raw compression}}
  = \mbox{\smChPerTB} .
\]
Numerator and denominator come from the same estimate, so the rate compares
like with like. One survey year in this model is \smNights{} nights $\times$
\smVisits{} visits per night, which is \smYrScale{} times the estimate's visit
list and \smInputTB{}\,TB of input. Each data release reprocesses every
survey year to date:
\[
  \text{core-hours} = \text{input TB per survey year} \times
  \text{survey years to date} \times \text{core-hours per input TB}.
\]
The first column is DP2 as actually run: \smDPtwoNodeDays{} node-days, all on
Milano nodes (DP2 performance metrics, June 2026).

\smpara{Peak cores.}
Peak concurrent cores spread each campaign over a \smWindowDays{}-day
processing window, inflated by an idle and inter-stage multiplier of \smIdle{}.
The multiplier is \emph{assumed}: the measured ratio of task wall time to task
CPU time is about 1.14, but time between stages has not been measured, and
1.3--1.5 has been suggested. A safety margin of \smSafety{} is shown
separately and is not included in peak demand.

\subsection{Fleet}

The USDF batch fleet has \smBatchNodes{} nodes: \smMilano{} Milano and
\smTorino{} Torino, each with \smCoresNode{} cores. Rubin pipelines run about
\smSpeedup{} times faster per core on Torino than on Milano, across all tasks.
Because the estimates are measured on Milano, capacity is expressed in
Milano-equivalent cores:
\[
  (\smMilano{} + \smTorino{} \times \smSpeedup{}) \times \smCoresNode{} = \mbox{\smMilanoEq{}} .
\]
The speedup was measured on one generation of tasks and should
be re-measured at each release.

\subsection{Result}

In year 2, with one survey year of data, the peak is \smDRonePeakCores{}
Milano-equivalent cores, or \smDRoneUtil{} of today's fleet. Every campaign
reprocesses the whole survey to date, so demand rises roughly linearly. With
all processing at the USDF, year 10 (nine survey years) needs \smDRnineUtil{}
of today's fleet (\tabref{tab:smCompute}). Sharing the
work with partner facilities changes this materially (\secref{sec:upsplit}).

\smpara{Cross-check against DP2.}
DP2 used about 32 node-days per thousand visits. The \vthirty{} DR1
estimate implies about 58, and the superseded \wfortyone{} estimate about 30.
This is not a contradiction: coadd-stage tasks scale with the number of
visits per patch, which is far higher in a full survey year than in DP2, so
the cost per visit should rise. The DRP team should still confirm
\vthirty{} as the planning basis, because the choice is worth roughly a
factor of two in DRP compute.

\begin{landscape}
% GENERATED by generate_model.py from lsst/rubin-sizing-model @ v2026.09.2. Do not edit by hand.
% Regenerate: git clone --branch v2026.09.2 https://github.com/lsst/rubin-sizing-model ; then, inside the clone, uv run generate_model.py --emit-tex <this directory>

\tiny \begin{longtable}{|p{0.19\textwidth}|l|r|r|r|r|r|r|r|r|r|r|}
\caption{DRP compute with all processing at the USDF, in core-hours, node-days and peak concurrent cores. The first column is DP2 as run. \label{tab:smCompute}}\\
\hline
\textbf{Metric} & \textbf{Unit} & \textbf{DP2 (actual)} & \textbf{DR1} & \textbf{DR2} & \textbf{DR3} & \textbf{DR4} & \textbf{DR5} & \textbf{DR6} & \textbf{DR7} & \textbf{DR8} & \textbf{DR9} \\
 &  & FY2026 & FY2027 & FY2028 & FY2029 & FY2030 & FY2031 & FY2032 & FY2033 & FY2034 & FY2035 \\
\hline\endfirsthead
\hline
\textbf{Metric} & \textbf{Unit} & \textbf{DP2 (actual)} & \textbf{DR1} & \textbf{DR2} & \textbf{DR3} & \textbf{DR4} & \textbf{DR5} & \textbf{DR6} & \textbf{DR7} & \textbf{DR8} & \textbf{DR9} \\
 &  & FY2026 & FY2027 & FY2028 & FY2029 & FY2030 & FY2031 & FY2032 & FY2033 & FY2034 & FY2035 \\
\hline\endhead
Survey years of input processed & years & 0.14 & 1.00 & 2.00 & 3.00 & 4.00 & 5.00 & 6.00 & 7.00 & 8.00 & 9.00 \\ \hline
Cumulative input data & TB & 114 & 798 & 1,596 & 2,395 & 3,193 & 3,991 & 4,789 & 5,588 & 6,386 & 7,184 \\ \hline
DRP work — core-hours & core-hrs & 2,793,600 & 35,481,751 & 70,963,503 & 106,445,254 & 141,927,006 & 177,408,757 & 212,890,508 & 248,372,260 & 283,854,011 & 319,335,762 \\ \hline
DRP work — node-days (at 120 cores/node) & node-days & 970 & 12,320 & 24,640 & 36,960 & 49,280 & 61,600 & 73,920 & 86,240 & 98,560 & 110,880 \\ \hline
DRP work — peak concurrent cores (in window) & cores & 794 & 10,079 & 20,158 & 30,237 & 40,316 & 50,395 & 60,474 & 70,553 & 80,632 & 90,711 \\ \hline
Core-hours with safety margin & core-hrs & 4,190,400 & 53,222,627 & 106,445,254 & 159,667,881 & 212,890,508 & 266,113,135 & 319,335,762 & 372,558,390 & 425,781,017 & 479,003,644 \\ \hline
Batch nodes needed (peak) & nodes & 7 & 84 & 168 & 252 & 336 & 420 & 504 & 588 & 672 & 756 \\ \hline
USDF fleet, Milano-equivalent cores & cores & 21,600 & 21,600 & 21,600 & 21,600 & 21,600 & 21,600 & 21,600 & 21,600 & 21,600 & 21,600 \\ \hline
Fleet utilization (peak / Milano-equivalent) & \% & 4\% & 47\% & 93\% & 140\% & 187\% & 233\% & 280\% & 327\% & 373\% & 420\% \\ \hline
\end{longtable} \normalsize

\end{landscape}
